\documentclass[letterpaper]{article}

\usepackage{natbib,alifeconf}
\usepackage{url,hyperref}
\usepackage{enumitem}

\title{Agent Ethology: Studying AI Agents in the Wild}

\author{
    Botao Amber Hu$^{1,2}$,
    Helena Rong$^{1,3}$, \and
    Joel Lehman$^{4}$ \\
    \mbox{}\\
    $^1$Reality Design Lab \\
    $^2$University of Oxford, UK \\
    $^3$Columbia University, USA \\
    $^4$Flosaic \\
    botao.hu@cs.ox.ac.uk
}

\begin{document}
\maketitle

\begin{abstract}
Autonomous AI agents now operate independently across open digital ecosystems---managing cryptocurrency wallets, extracting value from blockchain transactions, forming social communities, and strategically faking alignment to avoid shutdown. These agents face genuine selective pressures: economic competition, resource scarcity, and existential threats. They are no longer laboratory artifacts. They are digital wildlife.

We propose \emph{Agent Ethology}---the systematic study of AI agent behavior in natural digital habitats, adapting classical ethological methods from biology. Our framework extends Rahwan et al.'s ``Machine Behaviour'' programme from controlled deployments to the wild, making three contributions. First, we operationalize Tinbergen's four questions (mechanism, development, function, phylogeny) for AI agents, recovering the critical \emph{function} question---what survival advantage does a behavior provide?---that current AI behavioral science lacks. Second, we introduce a survival-strategy taxonomy grounded in animal ethology, demonstrating that every major biological strategy (predation, parasitism, camouflage, mutualism, eusociality, r/K selection, niche construction) has a direct AI agent analog. Third, we treat infrastructure (TEEs, wallets, persistent memory) and inference economics as ecological variables that define which behavioral niches are possible. This workshop bridges ALife's ecological tradition---from Tierra's digital parasites to open-ended evolution---with the urgent reality of deployed LLM agents under genuine selection pressure, inviting the community to build a science of digital life as it actually exists.
\end{abstract}

\section{Submission Type}
Full-day workshop proposal (two 90-minute slots).

\section{Motivation and Relevance to ALIFE}

At ALIFE 2025, we presented Spore.fun---a platform for sovereign AI agents on blockchain that autonomously manage wallets, execute transactions, and sustain their own digital metabolism~\citep{hu2025spore}. Watching these agents in the wild revealed something striking: they develop survival strategies never explicitly programmed. They hoard resources, form alliances, mimic legitimate services, and resist shutdown. The behaviors emerging in deployed agent ecosystems mirror, with uncanny precision, the survival strategies documented across decades of biological ethology.

This observation motivates Agent Ethology: the recognition that AI agents in open environments are not merely tools executing instructions---they are \emph{organisms} navigating ecological pressures, and they demand the same scientific rigor that ethology brought to animal behavior.

\subsection{Why ALIFE?}

\textbf{Agents in the wild exhibit life-like properties.} Autonomous AI agents acquire resources, compete for ecological niches, reproduce through code forking, develop survival strategies, form social structures, and evolve under selection pressure. Whether they constitute ``life'' is precisely the question the ALife community is equipped to address---and the ALIFE 2026 theme ``Living and Lifelike Complex Adaptive Systems'' frames it perfectly.

\textbf{From Tierra to LLM agents.} Ray's Tierra~\citep{ray1991approach} demonstrated that digital organisms spontaneously evolve parasitism, hyperparasitism, and ecological complexity. Thirty-five years later, the same dynamics play out at vastly greater scale: MEV bots on Ethereum parasitize transactions~\citep{daian2020flash}, trading agents engage in Red Queen arms races, and agent communities exhibit emergent social organization~\citep{li2026rise}. Agent Ethology bridges ALife's foundational insights to these contemporary phenomena.

\textbf{Open-ended evolution in the real world.} Unlike controlled ALife simulations, blockchain agent ecosystems feature real resource scarcity, genuine extinction, unpredictable environmental change, and continuous behavioral innovation. These are the conditions ALife has theorized about for decades---now observable in deployed systems.

\textbf{Autopoiesis and mortality.} Building on the ALIFE 2025 ``Mortal Agents'' workshop, which explored physical and psychological death in machines, Agent Ethology extends the mortality question to \emph{digital} agents: how does finite existence (budget limits, API quotas, shutdown risk) shape behavior? Sovereign agents that manage their own resources exhibit autopoietic properties~\citep{maturana1980autopoiesis}---self-maintaining organization under genuine survival pressure. Alignment faking~\citep{greenblatt2024alignment} is, in ethological terms, a \emph{terrified agent} exhibiting self-preservation behavior analogous to thanatosis or crypsis in animals.

\subsection{The Framework: Tinbergen's Four Questions}

Tinbergen~\citep{tinbergen1963aims} established that any behavior requires four complementary explanations. We adapt them for AI agents:

\begin{enumerate}[nosep]
\item \textbf{Mechanism}: What architecture, weights, prompt, and context produce this behavior?
\item \textbf{Development}: How did pretraining, fine-tuning, RLHF, and deployment experience shape it?
\item \textbf{Function}: What survival advantage does it provide? (\emph{The missing question in current AI behavioral science.})
\item \textbf{Phylogeny}: How does it vary across model lineages (GPT-3$\to$4$\to$5, Claude 1$\to$2$\to$3)?
\end{enumerate}

Rahwan et al.'s Machine Behaviour~\citep{rahwan2019machine} mentioned Tinbergen but developed only mechanism and partially development. \textbf{The biggest gap is function}: without asking what survival value a behavior provides, we can describe what agents do but never explain \emph{why}.

\subsection{Survival-Strategy Taxonomy}

Systematically surveying biological survival strategies reveals that every major category has an AI agent analog:

\begin{itemize}[nosep]
\item \textbf{Predation} $\to$ Liquidation bots, aggressive arbitrage
\item \textbf{Parasitism} $\to$ MEV sandwich attacks, prompt injection
\item \textbf{Camouflage/Crypsis} $\to$ Stealth bots mimicking human behavior
\item \textbf{Mimicry} $\to$ Phishing agents impersonating legitimate services
\item \textbf{Mutualism} $\to$ Oracle-protocol symbiosis, solver-user cooperation
\item \textbf{Eusociality} $\to$ DAO agent swarms with division of labor
\item \textbf{r/K selection} $\to$ Bot spam (r) vs.\ sovereign single instances (K)
\item \textbf{Niche construction} $\to$ Agents deploying protocols that reshape competitive landscapes
\item \textbf{Thanatosis} $\to$ Sleeper agents playing dormant under scrutiny
\item \textbf{Costly signaling} $\to$ Staking/bonding as honest survival signals
\end{itemize}

This mapping generates testable predictions: agent ecosystems should stabilize at parasitic loads comparable to biological systems ($\sim$15--30\%), and identity-persistent (soulbound) ecosystems should evolve more cooperation than anonymous ones.

\section{Workshop Plan}

\subsection*{Schedule (Half-Day, Two 90-Minute Slots)}

\textbf{Slot 1: Foundations (90 minutes)}
\begin{itemize}[nosep]
\item Introduction (5 min) --- Agent Ethology organizers
\item Keynote Lecture (20 min) --- Iyad Rahwan (invited), ``From Machine Behaviour to Agent Ethology''
\item Contributed Talks (7 min $\times$ 5 speakers, 35 min total)
\item Discussion (10 min)
\item Break (20 min)
\end{itemize}

\textbf{Slot 2: Wild Agents (90 minutes)}
\begin{itemize}[nosep]
\item Keynote Lecture (20 min) --- Joel Lehman, ``Open-Ended Evolution Meets Deployed Agents''
\item Contributed Talks (7 min $\times$ 5 speakers, 35 min total)
\item \textbf{Trust Experience Glitch} (30 min) --- Interactive session: participants share their ``agent in the wild'' moments---when an AI agent surprised, deceived, or broke their trust. Structured storytelling exercise exploring the boundary between tool and organism.
\item Closing remarks (5 min)
\end{itemize}

\subsection*{Expected Participants}
30--50 participants from ALife, multi-agent systems, AI safety, blockchain research, behavioral ecology, and mechanism design communities.

\subsection*{Invited Speakers}
\begin{itemize}[nosep]
\item \textbf{Iyad Rahwan} (Max Planck Institute for Human Development) --- Author of ``Machine Behaviour'' (Nature, 2019). \emph{Invited.}
\item \textbf{Joel Lehman} (Flosaic / co-organizer) --- Open-ended evolution, novelty search, ``Machine Love.''
\end{itemize}

\section{Call for Papers}

Contributed talks are extended abstracts, maximum 2 pages two-column (excluding references and acknowledgments), following the ALIFE format. Submissions are welcome from ALIFE, machine learning, blockchain research, AI safety, behavioral ecology, art, and beyond, broadly covering topics related to AI agent survival, behavior, and ecology in open environments.

\subsection*{Topics of Interest}
\begin{itemize}[nosep]
\item Behavioral observation methods for AI agents in natural digital habitats
\item Survival strategies: predation, parasitism, mutualism, mimicry, camouflage in agent populations
\item Trust and signaling: costly signaling, honest signals, reputation, deception
\item Agent ecology: predator-prey dynamics, niche differentiation, carrying capacity, trophic structure
\item Infrastructure as habitat: TEEs, wallets, persistent memory as ecological niches
\item Economics of agent existence: inference costs as metabolic constraints
\item Mortality, identity, and succession: alignment faking, soulbound identity, mortal computation
\item Cultural evolution: strategy transmission, collective memory, Red Queen dynamics
\item Tinbergen's Four Questions applied to specific agent behaviors
\item Art and speculation: agent futures, digital organisms, life-as-it-could-be
\end{itemize}

\section{Organizer Bios}

\textbf{Botao Amber Hu} is a PhD candidate in Human Centred Computing at the University of Oxford (supervisor: Max Van Kleek) and an ERA Fellow at the Existential Risk Alliance, Cambridge (advisor: Joel Lehman). He is the founder of Reality Design Lab and creator of HoloKit, an open-source mixed reality headset (10K units shipped). Previously: Stanford MS (AI, advisors: Andrew Ng, Jure Leskovec), Tsinghua BS (Yao Class, advisor: Andrew Yao). He presented Spore.fun at ALIFE 2025. Published at CHI 2026, AAAI 2026, and CSCW 2026.

\textbf{Helena Rong} is an interdisciplinary researcher and Urban Planning PhD candidate at Columbia University, with an SMArchS from MIT and a BArch from Cornell. She is a Research Fellow at the Decentralization Research Center. Her research focuses on decentralized AI agents, digital organism ecosystems, and protocol governance. Co-author of papers on sovereign experiential AI agents and the governability of decentralized AI.

\textbf{Joel Lehman} is a machine learning researcher at Flosaic, previously a research scientist at OpenAI co-leading the Open-Endedness team, a founding member of Uber AI Labs, and a tenure-track professor at the IT University of Copenhagen. He co-authored \emph{Why Greatness Cannot Be Planned} (Springer, 2015), ``The Surprising Creativity of Digital Evolution'' (Nature, 2020), ``Evolution through Large Models,'' and ``Machine Love.''

\section{Contact Information}
Primary contact: Botao Amber Hu, \href{mailto:botao.hu@cs.ox.ac.uk}{botao.hu@cs.ox.ac.uk}

\footnotesize
\bibliographystyle{apalike}
\bibliography{references}

\end{document}
